\documentclass[10pt]{article}
\usepackage[margin=20mm]{geometry}
\usepackage{amsmath,amssymb,fontspec,unicode-math,xurl,hyperref}
\IfFontExistsTF{DejaVu Serif}{\setmainfont{DejaVu Serif}}{\setmainfont{Latin Modern Roman}}
\setmathfont{latinmodern-math.otf}
\setlength{\parindent}{0pt}\setlength{\parskip}{6pt}\sloppy
\begin{document}
Abstract: This review checks the published latent-diffusion method, its computational evaluation, and a pinned released implementation. The supported assessment is Major, scoped to the release’s default ESM2/transformer decoder-training route under nonidentity normalization: training and noiseless validation apply the inverse normalization twice, while generation applies it once. This invalidates the shared latent-to-token interface on which the trained decoder’s use in generation depends. Three further Minor findings concern inconsistent reported metric values, repeated-sequence weighting in pLDDT, and an optional compressed-length mask. Independent algebra, finite examples, source inspection, and source-page images establish the discrepancies. Historical experiment-to-source bindings and quantitative benchmark effects remain unresolved; no replacement scores are asserted. Exact-sampling, biological-function, and matched-resource interpretations are bounded explicitly below. Scientific adjudication remains pending.\par
Date: 25 October 2026\par
Title: Scientific review of DiMA\par
\# Identity, version, and inspected scope\par
The paper is Diffusion on Language Model Encodings for Protein Sequence Generation, by Viacheslav Meshchaninov, Pavel Strashnov, Andrey Shevtsov, Fedor Nikolaev, Nikita Ivanisenko, Olga Kardymon, and Dmitry Vetrov, ICML 2025, PMLR 267 [Retained latex: \textbackslash{}cite\{paper\}]. Assignment identifiers are 0219 / I25-43588. All page references below refer to the 42-page publisher PDF; its printed Arabic page numbers agree with these PDF indices. The embedded appendix occupies pp. 15–42. The publisher-linked PDF was acquired on 25 October 2026 and is identified by its actual SHA-256 in Appendix 6 (\#app:evidence).\par
The released source is , commit , dated 17 July 2025 [Retained latex: \textbackslash{}cite\{code\}]. Its acquisition manifest does not assert that this commit generated any paper table. All 68 acquired members retain their archive, Git-tree, byte, and hash mappings. The two omitted members are documented NumPy length-distribution payloads; neither they nor datasets, model weights, normalization statistics, generated sequences, or historical run logs were acquired in this review.\par
The original parent reviewer read substantive main-text pp. 1–9 and appendix pp. 16–42, checked the equations and numerical comparisons discussed here, and independently inspected the central training, sampling, normalization, decoder, and metric source paths. Source pages 6, 26, 27, 29, 37, 40, and 42 were also inspected as rendered images. Other figure content received the coverage recorded in the ledger, not a claim of raw-data reanalysis. Three native auxiliary agents independently checked mathematics, source correspondence, and alternative interpretations. Their roles and limitations are preserved separately. No author code was executed.\par
\# Method and assumptions needed for the review\par
Equation reference: eq:schedule. DiMA encodes an amino-acid sequence into pretrained protein-language-model representations, normalizes each coordinate, learns a continuous denoiser, and decodes generated latent vectors. Writing the paper’s variance coefficient as \(A(t)\) to distinguish it from its square root, \[A(t)=\frac{1}{1+d^2\tan^2(\pi t/2)},\quad d=10,\qquad
 z_t=\sqrt{A(t)}z_0+\sqrt{1-A(t)}\epsilon,\quad \epsilon\sim N(0,I).
 \] Equation (2), p. 3, is a squared reconstruction objective with a detached self-conditioning prediction. At generation time the length is sampled separately, so the method need not learn the complete length distribution through padding tokens. The paper explicitly reports a quality–diversity trade-off, rather than simultaneous dominance on every metric.\par
The pLDDT evaluation in Appendix B.1, p. 16, first describes an average over residues within a predicted structure. Section 3.1, p. 4, describes evaluation over 2,048 generated sequences. These are model-based structural-confidence measurements. They do not, by themselves, certify experimentally measured folding or function. The authors acknowledge intrinsically disordered regions as a reason for low confidence; this caveat is retained rather than counted as an error.\par
The following findings distinguish a contradiction in published numbers, a deterministic source-level estimator change, and mismatches in released training/inference and compressed-length interfaces. Only the first directly concerns the published numerical presentation. The source-level findings require historical source/run bindings before their quantitative effect on that presentation can be established.\par
\# Supported problems\par
Problem 1 (Cluster-density definitions and reported values do not agree).  Appendix B.2, p. 17, defines positive clusters as co-clusters containing both generated and dataset sequences and negative clusters as co-clusters containing only generated sequences. It defines PCD and NCD as their respective counts divided by the total number of sequences. Tables 5 and 6, pp. 26–27, report a reference row with \(\mathrm{PCD}=0.990\) and \(\mathrm{NCD}=0.304\). Table 5 reports \(0.214\) and \(0.945\) for flow matching.\par
Complete check. Let \(G\) be the number of generated sequence instances and let \(P,N\) count the two kinds of clusters in the same clustering. Each counted cluster contains at least one generated instance, and clusters are disjoint. Thus \(P+N\leq G\). If the common denominator \(M\) is the number of generated instances, or the number of all generated-plus-reference instances, then \(M\geq G\), and \[\mathrm{PCD}+\mathrm{NCD}=\frac{P+N}{M}\leq 1.\] This argument does not require any monotonicity property of the clustering algorithm or an assumption that two thresholded clusterings are nested. The printed sums are \(1.294\) and \(1.159\). Even allowing independent rounding errors of \(0.0005\) in both entries, the sums are at least \(1.293\) and \(1.158\). They cannot represent the stated common-denominator quantities. ◻\par
Counterevidence and bounded impact. Separately normalized conditional densities, a transposed row, or an unstated denominator could explain the numbers, but each would require a corrected definition or table. The available repository lacks the co-clustering evaluation driver and underlying cluster counts, so the review does not choose replacement numbers. The contradiction prevents literal quantitative interpretation of these particular reference/ablation entries; it does not establish that the headline pLDDT or FD comparisons are wrong.\par
There is a separate, direct cross-table conflict in the ESMc-300M result: Table 3 and its discussion on p. 6 report \(\mathrm{CD}_{0.5}=0.963\), while Table 7 on p. 29 reports \(0.863\) for the corresponding AFDBv4-90 comparison. The difference is \(0.100\), not rounding. Both values remain below ESM2-650M’s \(0.986\), so the direction of that local comparison survives. The record must identify which value and run are authoritative. Likewise, the RFDiffusion unique success rate is described as \(0.06\) on pp. 8 and 36 but appears as \(0.07\), with an external-analysis footnote, in Table 13; either is below DiMA’s reported \(0.10\). These are Minor reporting defects, with exact corrections unverified. Parent arithmetic is in ; source-image anchors are listed in Appendix 6 (\#app:evidence).\par
Problem 2 (Released pLDDT aggregation discards within-rank multiplicity).  The inspected release’s , lines 31–40, stores a predicted score under the protein-sequence string: \(\texttt{result[protein] = score}\). The active caller in , lines 24–27, averages that dictionary’s values. Repeated occurrences on one rank overwrite their earlier entries.\par
Complete check. Assume a deterministic quality score \(q(u)\in[0,100]\) for each distinct sequence \(u\), and let \(m_u\) be its multiplicity among generated draws. The ordinary mean over draws is \[\bar q_{\rm draws}=\frac{\sum_u m_uq(u)}{\sum_u m_u},
 \qquad
 \bar q_{\rm unique}=\frac{1}{|\{u:m_u>0\}|}\sum_{u:m_u>0}q(u).\] The code computes the second quantity within each rank. Take four valid score-bearing draws \(A,A,A,B\), with hypothetical scores \(q(A)=90\), \(q(B)=30\). Then the generated-draw mean is \(75\), while the one-rank dictionary mean is \(60\). This example uses admissible score values, assumes no score stochasticity, and executes no folding model. If split into equally sized rank lists \((A,A)\) and \((A,B)\), the equal average of rank-local unique means is \((90+60)/2=75\). Thus even with equal shard lengths, the implemented result can depend on the partition. The independently authored finite calculation reproduces these values.\par
Counterevidence and impact. Quality of distinct outputs is a legitimate estimand when explicitly defined, especially alongside a diversity measure. It is not the same as average quality per generated attempt. A dictionary may also serve as a computation cache, but the final average must then restore each key’s multiplicity if the target is the draw mean. When all predictions are distinct, the present implementation agrees with the draw mean on equal-sized shards. The paper explicitly discusses duplicate generation in some baselines, so this distinction is relevant to comparative evaluation; it does not prove that this release evaluated those historical baselines, or establish the direction of any historical ranking change.\par
This is a Minor released-evaluator defect in its correspondence to the stated sample evaluation. Its effect on published core comparisons is Needs Evidence. A precise local repair is to retain a score per sample index, or cache by sequence while averaging with occurrence counts, and reduce sums/counts across ranks. Algebra verifies that this restores the intended finite-sample mean. No corrected protein score is claimed.\par
The same file averages PDB atom B-factors, whereas p. 16 describes residue averaging. These agree if residue weights are equal or appropriate aggregation is used; otherwise atom counts weight residues differently. The exact dependency/PDB convention was not independently run here, so that additional correspondence remains unresolved. A separate unequal-shard indexing issue exists when sample count is not divisible by world size, but the default counts 512/2,048 are divisible by the usual configured multi-GPU counts; it is not used to infer damage to reported runs.\par
Problem 3 (Decoder training and generation apply different inverse-normalization maps).  The method’s coordinate normalization and token reconstruction are described on pp. 3–4. The release’s documented decoder-training route calls . In its `loss\_step`, line 26 denormalizes the latent, then line 29 calls `batch\_get\_logits`. The ESM2 implementation, , lines 89–95, denormalizes again. In contrast, `batch\_decode`, lines 74–86, used by generation, denormalizes once. The training entry point enables the normalizer.\par
Core dependency. Figure 1, p. 2, identifies the decoder as one of the method’s three core components; the token-reconstruction paragraph, p. 3, describes its finetuning, and Table 5, p. 26, includes the finetuning ablation. The release’s , lines 3–4, selects ESM2 with the transformer decoder. , line 13, configures three new layers. , lines 37–47, constructs this decoder and loads its checkpoint when present; , lines 16–17, constructs the transformer and output projection from configuration. The documented Stage 2 command in , lines 128–142, trains that decoder through the affected route. Finally, , lines 492–532, sends generated normalized latents through `pred\_logits` and `batch\_decode`, using the single inverse map. The discrepancy therefore crosses the default learned decoder’s training, validation, and deployment boundary.\par
Equation reference: eq:double. Complete check. For one nondegenerate coordinate let the stored mean and standard deviation be \(\mu\) and \(s>0\), with \(s\geq 10^{-5}\) so the small-variance guard is inactive. Put \[N(x)=(x-\mu)/s,\qquad D(z)=sz+\mu.\] Training therefore presents \(D(D(z_t))\) to the decoder; generation presents \(D(z)\). Their difference at the same normalized input is \[D(D(z))-D(z)=s(s-1)z+s\mu.
 \] Equality for every \(z\) requires \(s=1,\mu=0\). The discrepancy also occurs in the noiseless decoder-validation branch, which bypasses corruption but not either denormalization.\par
For a valid finite interface example, take equally frequent raw encodings \(-1\) and \(3\). Their population mean is \(1\), variance \(4\), and normalized encodings are \(-1,+1\). The decoder’s training inputs are \(-1,7\), whereas the generation inputs are \(-1,3\). A two-class decoder with logits \(0,u-5\) correctly classifies targets \(0,1\) at the training inputs but assigns both generation inputs to class \(0\). Scaling both logits’ difference makes the training classification arbitrarily confident without repairing the inference mismatch. This establishes that fitting the double-transformed training inputs need not make the single-transformed inference path correct. It is an illustrative interface counterexample, not a claim that ESM2 produces these embeddings or that the released optimizer learns this classifier.\par
Counterevidence, repair, and impact. The discrepancy vanishes for identity statistics, at particular fixed points, or for a decoder insensitive to the transformed directions. The README also permits skipping decoder training with the pretrained `lm\_head` (line 132); that alternative does not establish correctness of the default newly trained transformer path. Other encoders or historically different code are outside this finding’s demonstrated dependency. Actual normalization tensors and trained decoder weights are absent, so the example cannot estimate historical accuracy loss. A model might tolerate some mismatch. These possibilities bound the affected configurations and numerical impact; they do not make the default training and deployment interfaces equal.\par
Removing the outer denormalization from `loss\_step`, while retaining the existing denormalization inside `batch\_get\_logits`, restores the single map \(D(z_t)\) algebraically for every admitted \(\mu,s,z_t\). This is an interface repair only: it does not validate old checkpoints or prove the paper’s reported finetuning improvement without retraining/evaluation. The supported assessment is Major for the pinned release’s default ESM2/transformer decoder-training route under nonidentity statistics. A core learned component is optimized and validated on a different latent-to-token interface from the one used in generation; the finite example proves that successful fitting and noiseless validation need not transfer through that interface. Quantitative historical effects and the paper’s reported finetuning gain remain Needs Evidence. The scope is this demonstrated core training/deployment inconsistency, without a claim that every encoder, checkpoint, or published protein-generation result fails.\par
Problem 4 (The optional shorten-two mask applies a second length reduction).  The CHEAP shorten-two component is evaluated in Sections 3.4.2 and E.6. In the pinned release, , lines 74–86, first compresses the batch maximum length by two, then clips each uncompressed sequence length to that already compressed maximum before halving it again.\par
Complete check. Take the valid interface inputs \(\texttt{lens}=[254]\), \(\texttt{max\_sequence\_len}=254\), and \(\texttt{shorten\_factor}=2\). The allocated latent width is \(\lceil254/2\rceil=127\), but the per-sequence mask receives only \[\left\lceil\frac{\min(254,127)}2\right\rceil=64\] active positions, rather than 127. The decoding path at lines 70–71 explicitly truncates the sequence to twice the mask sum, so the output has at most 128 residues, not the requested 254. The length sampler returns raw sequence lengths and supplies no prior factor-two conversion. This calculation depends only on the explicit mask and output-slice arithmetic, not on the external CHEAP decoder’s behavior.\par
Counterevidence and impact. The supplied default and supported pretrained configuration use shorten-one, where this particular branch is inactive. No historical shorten-two configuration or generated-length records bind this code to the corresponding paper results. The supported judgment is therefore Minor for the optional released path, with historical impact unresolved. A precise repair is to clamp each raw length to the raw maximum before applying the single compression: \(m_i=\lceil\min(\ell_i,L_{\max})/2\rceil\), bounded by the allocated compressed width. For odd lengths the decoder must also retain the requested raw length if exact truncation is intended. This restores the mask count; it does not validate external decoder integration or old experiments. The parent calculation is recorded in .\par
\# Checks that do not establish additional core failures\par
\#\# Diffusion coefficients and self-conditioning\par
For a variance-preserving forward SDE with drift \(-\beta(t)z/2\) and diffusion \(\sqrt{\beta(t)}\), consistency with [eq:schedule] (\#eq:schedule) requires \[\beta(t)=-\frac{d}{dt}\log A(t)
 =\frac{\pi d^2\tan(\pi t/2)\bigl(1+\tan^2(\pi t/2)\bigr)}
 {1+d^2\tan^2(\pi t/2)}.\] The release’s , lines 36–52, uses this expression but caps its time argument at \(0.95\), while its marginal coefficients use uncapped time. and the configured Euler solver actively use the capped coefficient. Thus those two functions do not define exactly the same forward SDE over \(0.95<t<1\).\par
The discrepancy is quantified rather than assumed catastrophic. At \(d=10\), \(A(0.95)=6.1935750\times10^{-5}\), \(\beta(0.95)=40.1624805\), and \(\beta(0.99)=200.0324088\). The capped forward process would have \[A_{\rm cap}(1)=A(0.95)\exp[-0.05\beta(0.95)]
 =8.3142718\times10^{-6},\] whereas the uncapped formula tends to zero. For standardized Rademacher data \(z_0=\pm1\) with equal probabilities, both processes retain mean zero and variance one. Their endpoint fourth cumulants are respectively \(-2A_{\rm cap}(1)^2=-1.3825423\times10^{-10}\) and zero. This is a legitimate consistency witness and also strong counterevidence to claiming a large effect from this example. It is neither a general high-dimensional error bound nor an estimate of the reverse sampler’s actual protein error.\par
The default sampling grid ends at \(t=0.05\) and returns the final solver mean, not its predicted clean \(z_0\); here \(A(0.05)=0.6175142\) and the forward noise standard deviation is \(0.6184543\). Exact clean-latent sampling cannot simply be presumed from this interface. However, Appendix C.5, p. 23, explicitly motivates robustness at \(t=0.05\), reporting near-lossless CHEAP reconstruction. The paper offers an empirical schedule and no exact finite-step sampling theorem. Numerical endpoint regularization and intentional early decoding are therefore credible interpretations. Their effects across encoders remain Needs Evidence, not an additional demonstrated Major problem.\par
For an exact denoiser, the source conversion \(\widehat{\mathrm{score}}=-(z_t-\sqrt{A(t)}\,\widehat z_0)/(1-A(t))\) has the expected sign and scaling, and the Euler reverse-step noise uses \(\sqrt{-dt}\) for decreasing time. The current training code computes both detached self-conditioned and zero-conditioned losses on each batch. Their sum has twice the expected 50/50 mixture objective in paper Eq. (2). The missing random branch coin alone is not an invalid objective; proportional objectives do not imply identical finite optimizer trajectories under clipping and other settings. The tan-10 rationale is presented as a heuristic, not a theorem that every representation’s loss must be linear in time.\par
\#\# Empirical comparisons and conditional tasks\par
The stated Gaussian Fréchet formula and the biased equal-size MMD formula in Appendix B.3 pass the inspected algebraic checks. The released masked-token pseudoperplexity excludes special tokens from its average and follows the stated conditional-log-probability construction. These are scoped checks of formulas and inspected paths, not empirical reproductions.\par
The 35M comparison counts the denoiser, while DiMA benefits from a pretrained encoder during training. Removing the encoder at deployment can support a smaller generation model; it does not imply equal total pretraining resources. The phrase that all models are trained from scratch in Section 3.3/Table 6 needs this qualification. The paper discloses pretrained encoders, and the no-encoder ablation is useful counterevidence against concealment. A causal claim that only the diffusion paradigm explains the performance gap would require a matched-representation/resource comparison not supplied here.\par
For motif scaffolding, independently summing the 24 DiMA entries in Table 13 yields 19 positive success-rate tasks and mean success rate \(0.1658333\), agreeing with 19/24 and the printed \(0.17\). This check supports the numerical headline. It does not reproduce structures or unique-success clustering. Table 13 also includes a sequence-mode ESM3 result with 21 solved tasks, so the main-text comparison against named EvoDiff/DPLM/DPLM2 baselines should not be broadened to every sequence-based model.\par
For fold conditioning, Appendix E.9 explicitly defines the mean as the best TM-score among ten attempts per target, and acknowledges that DiMA receives richer conditioning than RFDiffusion’s secondary-structure/block-adjacency information. The reported \(0.93\) is not an unconditional per-attempt mean or a comparison at equal information. Nonzero RMSD alone cannot exclude copying followed by imperfect reconstruction, but Figure 5’s selected low sequence identities provide counterevidence to a blanket exact-sequence-copy allegation. No such allegation is established.\par
The family-task evaluation filters to correctly annotated family members before reporting quality/diversity, while reporting fidelity separately. That conditional estimand is disclosed on p. 31. The infilling comparison uses ten attempts on 180 targets; its 42.2\% versus 40.0\% point estimates do not establish a statistically resolved advantage without per-target outcomes. Biological annotation, predicted disorder, and predicted structures support computational plausibility; no wet-lab functional validation was established or represented as having been performed.\par
\# Dependence, severity, and remaining work\par
The supported overall impact is Major, scoped to the pinned release’s default ESM2/transformer decoder-training and generation route under nonidentity normalization. Problem Problem 3 (\#prob:decoder) establishes an invalid input contract across a core component’s learning, validation, and use. The metric-definition, multiplicity, and optional-mask findings remain Minor for the bounded reasons stated individually. Historical source bindings and corrected benchmark values remain unresolved. This does not lower the demonstrated core-interface error to Minor: severity follows its operative dependency, while the extent of historical numerical harm remains a separate question. The judgment does not depend on appendix placement, repair difficulty, a severity quota, or whole-paper collapse.\par
No later corrected PDF or author-adopted repair was established from the frozen inputs. The optional OpenReview metadata and separate-supplement endpoint had returned HTTP 403 during acquisition; that does not establish that no additional material exists. The complete available source is sufficient for the static checks but omits conditional-task drivers, several reported metric pipelines, and historical run artifacts. Current defaults also differ from the narrative: a three-layer transformer decoder is configured instead of the described simple decoder, and pLDDT/pseudoperplexity use 512 samples while FD/MMD use 2,048. These are version-correspondence limitations, not proof of the published settings.\par
Item | Judgment | Repair verification | Adoption / remaining work\par
Metric definitions and tables | Minor | Constraint proved; replacement numbers unknown | Not requested; raw counts/run identity\par
pLDDT multiplicity | Minor, release | Sum/count repair proved for stated mean | Not requested; historical sequences/scores\par
Decoder normalization | Major, release | Single-map interface repair proved | Not requested; stats, checkpoints, evaluation\par
Compressed-length mask | Minor, optional | Correct count proved | Not requested; historical shorten-two runs\par
Core quantitative impact | Needs Evidence | No corrected benchmark asserted | Historical bindings and affected reruns\par
All scientific adjudications remain `pending`; `human\_reviews` is empty. Repair verification and author adoption are distinct states.\par
9 V. Meshchaninov et al., Diffusion on Language Model Encodings for Protein Sequence Generation, Proceedings of ICML, PMLR 267 (2025). Publisher record: https://proceedings.mlr.press/v267/meshchaninov25a.html. Reviewed publisher PDF: https://raw.githubusercontent.com/mlresearch/v267/main/assets/meshchaninov25a/meshchaninov25a.pdf. The appendix is embedded in the same frozen PDF. V. Meshchaninov et al., DiMA source repository, https://github.com/MeshchaninovViacheslav/DiMA/tree/18f4a2e67988efe1cc5593fcaec1ece8f09fdac0, 17 July 2025. Consulted through the preserved local archive and mapped source members, without author-code execution.\par
\# Evidence, calculations, and provenance\par
The original review was produced in\par
> .\par
The self-contained source/PDF pair is and . Supporting records reside in . The input PDF is ; source member paths above are relative to .\par
Actually verified immutable inputs.\par
- Publisher PDF SHA-256: .\par
- Source-manifest SHA-256: .\par
- Preserved selective ZIP SHA-256: .\par
- Member-map SHA-256: .\par
The source check found zero mapped byte/hash mismatches and a maximum checked source-path length of 142 characters. The member map preserves meaningful repository hierarchy and the two excluded payload entries. No source identity, version, or bytes were changed.\par
Independent checks and evidence anchors. The parent-authored uses only the Python standard library; records actual outputs for the PCD/NCD inequality, ESMc discrepancy, multiplicity example, double-denormalization and compressed-length examples, tan-10 coefficient calculation, and Table 13 summation. It imports no author modules. These are exact finite arithmetic and analytic diagnostics, not a run of the trained DiMA model. The image evidence is , , , , , , and ; it was used to exclude extraction ambiguity in the relevant table entries and inspect selected biological figures. The auxiliary records are , , and . They are AI cross-checks, not human reviews.\par
The original PDF was extracted with the existing `pdftotext -layout` and selected pages rendered with `pdftoppm`. An initial page-read attempt failed because had not yet been created; extraction then completed successfully and the reads were repeated. No failed read is counted as access. A report-draft JavaScript delimiter error occurred before any TeX file was written; the draft was corrected before compilation. There were no downloads, installs, author-code runs, dataset/model acquisitions, or experimental measurements.\par
Compilation and visual verification. The supported isolated runtime has the existing `/usr/bin/pdflatex`, pdfTeX 3.141592653-2.6-1.40.25 (TeX Live 2023/Debian). The report uses the frozen mathematical-journal template’s `amsart`, 11-point, single-column structure and locally available packages. The parent executes two passes with `-no-shell-escape -interaction=nonstopmode -halt-on-error -file-line-error`, renders the compiled PDF, and inspects every report page. , compilation logs, and bind the actual final source, compiled PDF, page count, and visual result; records their actual byte hashes and paths. These receipts, rather than a filename alone, establish the final artifact correspondence.\par
Roles and decisions. The user’s selected conversation label is 6 Pro; the exposed system model label is GPT-6 Astra Pro. The parent performed independent scientific integration, source checks, calculations, compilation, repairs if needed, and visual inspection. The native auxiliaries , , and performed the limited roles recorded above. Their separate model identifiers, an independently observed UI label, the conversation URL, and accessible Space identifiers were not exposed and are recorded as null, not invented. Those role facts describe the original review. The focused post-batch acceptance checked the default decoder dependency and corrected its severity; actual acceptance roles, final byte hashes, compilation, and visual checks are recorded in the existing . No cloud copy, author adoption, or human attestation is asserted. The frozen assignment token is . Hashes of the actually read packet, prompt, template, and review contract are preserved in . The coverage ledger distinguishes inspected, partial, unavailable, and unperformed work.\par
\end{document}
